% ============================================================================
%  TEMPLATE D'AFFICHE SCIENTIFIQUE DATAIA (FORMAT VERTICAL / PORTRAIT — A0)
%  Institut DataIA, Université Paris-Saclay
%  Compilation : XeLaTeX (obligatoire, pour la police Raleway)
% ============================================================================
%
%  COMMENT PERSONNALISER CE TEMPLATE ?
%  Tout ce dont vous avez besoin se trouve dans le bloc
%  "PARAMÈTRES À PERSONNALISER" juste en dessous : logo, couleurs,
%  informations de l'auteur, contact. Le reste du fichier n'a normalement
%  pas besoin d'être modifié.
%
%  Le centre de l'affiche contient 3 blocs de contenu prêts à l'emploi
%  (voir section "CONTENU DE L'AFFICHE" tout en bas) : à vous de les
%  déplacer, redimensionner ou dupliquer selon votre contenu scientifique.
% ============================================================================

\documentclass[a0paper]{article}

\usepackage{fontspec}
\usepackage{polyglossia}
\setmainlanguage{french}
\usepackage[paperwidth=84.1cm, paperheight=118.9cm, margin=0cm]{geometry}
\usepackage{graphicx}
\usepackage{grffile} % permet les noms de fichiers avec espaces/parenthèses (ex. "Image_défault (1).jpg")
\usepackage{tikz}
\usepackage{enumitem} % pour personnaliser l'espacement des listes à puces

\pagestyle{empty}
\setlength{\parindent}{0pt}

% ----------------------------------------------------------------------------
% PARAMÈTRES À PERSONNALISER
% ----------------------------------------------------------------------------

% --- Polices --------------------------------------------------------------
% Les fichiers .ttf de Raleway doivent se trouver dans le dossier Fonts/
% (déjà inclus dans le projet Overleaf, réutilisé depuis le template de
% slides). Si Raleway n'est pas disponible, remplacez "Raleway" par le
% nom d'une police installée sur Overleaf.
\setmainfont{Raleway}[
    Path = Fonts/,
    Extension = .ttf,
    UprightFont = *-Regular,
    BoldFont = *-Bold,
    ItalicFont = *-Regular,
]
\newfontfamily\ralewaylight{Raleway}[Path = Fonts/, Extension = .ttf, UprightFont = *-Light]
\newfontfamily\ralewaymedium{Raleway}[Path = Fonts/, Extension = .ttf, UprightFont = *-Medium]
\newfontfamily\ralewaysemibold{Raleway}[Path = Fonts/, Extension = .ttf, UprightFont = *-SemiBold]

% --- Couleurs officielles DataIA -------------------------------------------
\definecolor{dataiaBleuFonce}{HTML}{253573}   % couleur principale (bandeaux)
\definecolor{dataiaBleuMarine}{HTML}{3997D2}  % accent
\definecolor{dataiaGris}{HTML}{4A4949}        % texte secondaire
\definecolor{dataiaVert}{HTML}{B4CC0D}
\definecolor{dataiaRose}{HTML}{DA398D}
\definecolor{dataiaOrange}{HTML}{ED700E}
\definecolor{dataiaBlanc}{HTML}{FFFFFF}

% --- Logo et badge -----------------------------------------------------
% Chemins vers vos visuels, dans le dossier Image/ (les mêmes fichiers
% que ceux utilisés pour le template de slides).
\newcommand{\logoDataIA}{Image/logo-dataia-couleur.png}
\newcommand{\badgeSupportedBy}{Image/badge-supported-by.png}

% --- Informations de l'auteur (bandeau du haut) ------------------------
\newcommand{\posterNom}{[Prénom] [Nom]}
\newcommand{\posterRole}{[Rôle] -- [Affiliation]}
\newcommand{\posterEcole}{[Faculté / École]}

% --- Contact (bandeau du bas) ------------------------------------------
\newcommand{\posterEmail}{prenom.nom@universite-paris-saclay.fr}

% --- Dimensions des bandeaux ---------------------------------------------
% Ajustez ces deux valeurs si vous voulez un bandeau du haut plus haut/bas,
% ou un bandeau du bas plus ou moins fin.
\newlength{\hauteurBandeauHaut}
\setlength{\hauteurBandeauHaut}{12cm}
\newlength{\hauteurBandeauBas}
\setlength{\hauteurBandeauBas}{4cm}

% ----------------------------------------------------------------------------
% CONFIGURATION (ne pas modifier sauf besoin avancé)
% ----------------------------------------------------------------------------

% Largeur/hauteur de la page, pratiques pour les tikzpicture en overlay
\newlength{\posterlargeur}
\newlength{\posterhauteur}
\setlength{\posterlargeur}{\the\paperwidth}
\setlength{\posterhauteur}{\the\paperheight}

\begin{document}

% ============================================================================
% BANDEAU DU HAUT — fond bleu DataIA, infos auteur à gauche, logo à droite
% ============================================================================
\begin{tikzpicture}[remember picture, overlay]
    % Fond bleu plein du bandeau du haut
    \fill[dataiaBleuFonce] (current page.north west) rectangle
        ([yshift=-\hauteurBandeauHaut]current page.north east);

    % Nom / Prénom
    \node[anchor=north west, text=white, font=\ralewaysemibold, xshift=2.5cm, yshift=-3cm]
        at (current page.north west) {\fontsize{48}{58}\selectfont \posterNom};

    % Rôle -- Affiliation
    \node[anchor=north west, text=white, font=\ralewaymedium, xshift=2.5cm, yshift=-6cm]
        at (current page.north west) {\fontsize{34}{34}\selectfont \posterRole};

    % Faculté / École
    \node[anchor=north west, text=white, font=\ralewaylight, xshift=2.5cm, yshift=-7.8cm]
        at (current page.north west) {\fontsize{34}{34}\selectfont \posterEcole};

    % Logo DataIA, à droite du bandeau
    \node[anchor=north east, xshift=-2.5cm, yshift=-2.6cm]
        at (current page.north east) {\includegraphics[height=6cm]{\logoDataIA}};
\end{tikzpicture}

% ============================================================================
% BANDEAU DU BAS — fond bleu DataIA (fin), email à gauche, badge à droite
% ============================================================================
\begin{tikzpicture}[remember picture, overlay]
    % Fond bleu plein du bandeau du bas
    \fill[dataiaBleuFonce] (current page.south west) rectangle
        ([yshift=\hauteurBandeauBas]current page.south east);

    % Email de contact, à gauche
    \node[anchor=west, text=white, font=\ralewaymedium, xshift=2.5cm]
        at ([yshift=\hauteurBandeauBas/2]current page.south west) {\fontsize{30}{24}\selectfont \posterEmail};

    % Badge "Supported by the DataIA Institute", à droite
    \node[anchor=east, xshift=-2.5cm]
        at ([yshift=\hauteurBandeauBas/2]current page.south east) {\includegraphics[height=14cm]{\badgeSupportedBy}};
\end{tikzpicture}

% ============================================================================
% CONTENU DE L'AFFICHE
% ============================================================================
% L'espace entre les deux bandeaux est votre zone de contenu scientifique.
% 3 blocs prêts à l'emploi sont fournis ci-dessous, chacun dans son propre
% tikzpicture (ne jamais imbriquer deux tikzpicture l'un dans l'autre).
% Déplacez-les en changeant simplement les valeurs de xshift/yshift, ou
% dupliquez un bloc pour en ajouter un nouveau.

% --- Bloc 1 : titre + paragraphe — en haut à gauche -------------------------
\begin{tikzpicture}[remember picture, overlay]
    \node[anchor=north west, xshift=2.5cm, yshift=-\hauteurBandeauHaut-7cm,
          text width=0.42\posterlargeur, align=left]
        at (current page.north west) {%
            {\ralewaysemibold\fontsize{38}{44}\selectfont [Titre]}\\[0.8cm]
            {\ralewaylight\fontsize{26}{34}\selectfont
            Lorem ipsum dolor sit amet, consectetur adipiscing elit. Sed do
            eiusmod tempor incididunt ut labore et dolore magna aliqua. Ut
            enim ad minim veniam, quis nostrud exercitation ullamco laboris
            nisi ut aliquip ex ea commodo consequat. Duis aute irure dolor in
            reprehenderit in voluptate velit esse cillum dolore eu fugiat
            nulla pariatur. Excepteur sint occaecat cupidatat non proident,
            sunt in culpa qui officia deserunt mollit anim id est laborum.
            Curabitur pretium tincidunt lacus, sed rhoncus nisl porta a. Sed
            ac lacus at libero pretium fermentum. Aenean euismod, nunc eu
            fermentum tincidunt, nunc lacus vehicula lacus, eget aliquet
            lacus lacus eget lacus. Nulla facilisi. Vestibulum ante ipsum
            primis in faucibus orci luctus et ultrices posuere cubilia
            curae; Donec sit amet nisl at leo tincidunt tincidunt.}
        };
\end{tikzpicture}

% --- Bloc 2 : image + légende — en haut à droite -----------------------------
\begin{tikzpicture}[remember picture, overlay]
    \node[anchor=north east, xshift=-2.5cm, yshift=-\hauteurBandeauHaut-7cm]
        (imageBlocUn) at (current page.north east)
        {\includegraphics[width=0.42\posterlargeur]{Image/Image_défault (1).jpg}};
    \node[anchor=north east, xshift=0cm, yshift=-0.8cm,
          text width=0.42\posterlargeur, align=left]
        at (imageBlocUn.south east) {%
            {\ralewaylight\fontsize{24}{30}\selectfont
            Lorem ipsum dolor sit amet, consectetur adipiscing elit,
            sed do eiusmod tempor incididunt ut labore et dolore magna.}
        };
\end{tikzpicture}

% --- Bloc 3 : titre + liste à puces — au centre à gauche ---------------------
\begin{tikzpicture}[remember picture, overlay]
    \node[anchor=north west, xshift=2.5cm, yshift=-\hauteurBandeauHaut-25cm,
          text width=0.42\posterlargeur, align=left]
        at (current page.north west) {%
            {\ralewaysemibold\fontsize{38}{44}\selectfont [Titre]}\\[0.8cm]
            \begin{itemize}[leftmargin=1.2cm, itemsep=0.5cm, label=\textbullet]
                \item {\ralewaylight\fontsize{26}{32}\selectfont ...}
                \item {\ralewaylight\fontsize{26}{32}\selectfont ...}
                \item {\ralewaylight\fontsize{26}{32}\selectfont ...}
            \end{itemize}
        };
\end{tikzpicture}

% --- Bloc 4 : titre + paragraphe (2e bloc texte) — au centre à droite -------
\begin{tikzpicture}[remember picture, overlay]
    \node[anchor=north east, xshift=-2.5cm, yshift=-\hauteurBandeauHaut-37cm,
          text width=0.42\posterlargeur, align=left]
        at (current page.north east) {%
            {\ralewaysemibold\fontsize{38}{44}\selectfont [Titre]}\\[0.8cm]
            {\ralewaylight\fontsize{26}{34}\selectfont
            Lorem ipsum dolor sit amet, consectetur adipiscing elit. Sed do
            eiusmod tempor incididunt ut labore et dolore magna aliqua. Ut
            enim ad minim veniam, quis nostrud exercitation ullamco laboris
            nisi ut aliquip ex ea commodo consequat. Duis aute irure dolor in
            reprehenderit in voluptate velit esse cillum dolore eu fugiat
            nulla pariatur. Excepteur sint occaecat cupidatat non proident,
            sunt in culpa qui officia deserunt mollit anim id est laborum.
            Curabitur pretium tincidunt lacus, sed rhoncus nisl porta a. Sed
            ac lacus at libero pretium fermentum. Aenean euismod, nunc eu
            fermentum tincidunt, nunc lacus vehicula lacus, eget aliquet
            lacus lacus eget lacus. Nulla facilisi. Vestibulum ante ipsum
            primis in faucibus orci luctus et ultrices posuere cubilia
            curae; Donec sit amet nisl at leo tincidunt tincidunt.}
        };
\end{tikzpicture}

% --- Bloc 5 : image + légende (2e bloc image) — en bas à gauche ------------
\begin{tikzpicture}[remember picture, overlay]
    \node[anchor=north west, xshift=2.5cm, yshift=-\hauteurBandeauHaut-39cm]
        (imageBlocDeux) at (current page.north west)
        {\includegraphics[width=0.42\posterlargeur]{Image/Image_défault (1).jpg}};
    \node[anchor=north west, xshift=0cm, yshift=-0.8cm,
          text width=0.42\posterlargeur, align=left]
        at (imageBlocDeux.south west) {%
            {\ralewaylight\fontsize{24}{30}\selectfont
            Lorem ipsum dolor sit amet, consectetur adipiscing elit,
            sed do eiusmod tempor incididunt ut labore et dolore magna.}
        };
\end{tikzpicture}

\end{document}
```